\section{Introduction}
\label{sec:intro}

\subsection{The Problem of Vulnerability Reproduction}

The Common Vulnerabilities and Exposures (CVE) program assigns an identifier
and a short description to every publicly disclosed vulnerability, and the
number of entries grows by tens of thousands every year~\cite{cve,nvd}. A CVE
entry is a pointer, not an experiment: it names the affected product and
versions, and it may link to an advisory or a patch, but it does not tell the
reader how to bring the vulnerable software up, how to trigger the flaw, or
how to tell whether the attempt worked. Turning an entry into a working
reproduction is left to whoever needs it.

Many people need it. Reproductions are the ground truth on which much of
security research and practice rests:

\begin{itemize}
\item \textbf{Detection models.} Network and host intrusion detection systems,
signature generators, and machine-learning anomaly detectors must be trained
and evaluated on traffic and system activity that contains real, successful
attacks, alongside comparable benign activity. Long-lived public datasets are
either synthetic or stale, and evaluating detectors on them has been
repeatedly criticized~\cite{sommer2010outside,lippmann2000darpa,creech2013adfa}.
A reproduction that can be rerun at will produces fresh, labeled captures for
exactly the vulnerability of interest.

\item \textbf{Patch analysis.} Security patches are often incomplete:
large-scale studies show that a notable fraction of fixes must be revised
before a vulnerability is really closed~\cite{li2017patches}. The two Apache
\texttt{httpd} path-traversal vulnerabilities in our corpus illustrate this:
the fix for CVE-2021-41773 shipped in version 2.4.50 was itself bypassable,
which led to CVE-2021-42013. An executable reproduction turns ``does the
patch work?'' into a test that can be run against the patched build, and
against every later version.

\item \textbf{Automated security tools.} Fuzzers, program analyzers, exploit
generators, and, more recently, LLM-based agents are all evaluated on
vulnerabilities whose presence and exploitability must be established
beforehand~\cite{hazimeh2020magma,zhu2025cvebench,wang2026cybergym}. Without
reproducible environments, these evaluations are hard to compare and hard to
repeat.

\item \textbf{Risk assessment and education.} Whether a vulnerability is
exploitable in a given configuration, and what an exploitation looks like on
the wire and on the host, are questions that defenders answer most reliably
by running the attack. Reproducing attacks is also one of the most effective
ways to teach how they work.
\end{itemize}

Reproduction, however, is expensive. In the most extensive study of the
subject, a team of analysts spent 3,600 hours trying to reproduce 368
crowd-reported memory-corruption vulnerabilities, and found that reports
routinely omit the information needed to succeed~\cite{mu2018reproducibility}.

\subsection{The Challenge: No Standard Reproduction Environment}

The difficulty is not only in crafting an exploit. Most of the effort goes
into the environment: obtaining the exact vulnerable version, building it
with a compatible toolchain, installing its dependencies, configuring it into
an exploitable state, and seeding it with the data an attack needs. Every
reproduction effort solves these problems in its own way, and the solution is
rarely captured in a reusable form. Exploit databases publish exploits but
not the environment they target~\cite{exploitdb}. Collections of
containerized vulnerable applications, such as Vulhub~\cite{vulhub}, provide
environments but no uniform notion of what a successful exploitation is, and
no instrumentation. Benchmarks built for fuzzers or AI agents define
environments tailored to their own harnesses~\cite{mei2026arvo,zhu2025cvebench,wang2026cybergym,lee2025secbench}.
As a result, there is no standard way to \emph{describe} the environment in
which a vulnerability is reproduced, to \emph{verify} that the reproduction
succeeded, or to \emph{record} what happened during the attack. Results
produced by one group cannot easily be rerun by another, and reproductions
are lost when the original setup disappears.

\subsection{Our Solution: CVEX, the Executable CVE}

We propose \cvex, short for \emph{executable CVE}: a standard format that
augments a CVE entry with everything needed to reproduce it. A \cvex record
is a small, self-contained directory that describes (i)~the vulnerable
\emph{target}, built in a container; (ii)~a \emph{benign} workload that uses
the same service legitimately; (iii)~an \emph{attack} that exploits the
vulnerability; (iv)~a \emph{recorder} that captures network packets and
system calls; and (v)~a machine-checkable \emph{success criterion}. Running a
record produces a verdict and a dataset.

The central design choice of the third version of the format, \cvexthree, is
how success is defined. Following Capture-the-Flag (CTF) competitions, and
the verifiers produced by the \genie automated reproduction
framework~\cite{ullah2026cvegenie}, a record plants a secret \emph{flag} in
the target, at a location that only a successful exploitation can reach: a
root-owned file for a privilege escalation, a file outside the web root for a
path traversal, a database row for an SQL injection. The attack succeeds if
and only if it returns the flag. This makes verification objective and
uniform across vulnerability classes, and it guarantees that the attack half
of every dataset contains an exploitation that actually worked.

\subsection{Contributions}

This report makes two contributions.

\begin{enumerate}
\item \textbf{The \cvexthree format.} We define a uniform, Docker-based record
structure with a declarative manifest, three deployment topologies, and three
success criteria (planted flag, CTF-style verifier, and crash witness). We
define a run protocol that exercises every record in two isolated phases,
benign and attack, each against a pristine target, and records both as packet
captures and system-call traces. We provide a reference runner that builds a
record, runs it, verifies it, and writes a content-addressed dataset
(Section~\ref{sec:format}).

\item \textbf{A corpus of executable CVEs.} We release \NumRecords \cvexthree
records. \NumFunctional of them are fully functional: they have been built
and run by the reference runner, their exploitation has been verified, and
they ship with reference benign/attack captures. These include \NumMigratedFunctional
records migrated by hand from the previous generation of the format
and \NumGenieFunctional records converted from \genie reproductions, covering
\NumGenieCWEs CWE classes across \NumGenieProjects open-source projects. The
remaining records are scaffolds that are being functionalized
(Section~\ref{sec:dataset}).
\end{enumerate}

The rest of this report is organized as follows. Section~\ref{sec:background}
reviews CVEs, CWEs, existing approaches to reproduction, and the earlier
versions of \cvex. Section~\ref{sec:format} describes the \cvexthree format,
its runner, and the corpus. Section~\ref{sec:examples} walks through two
records in detail. Section~\ref{sec:future} discusses extending the format to
virtual machines, Section~\ref{sec:related} surveys related work, and
Section~\ref{sec:conclusion} concludes.
